% =====================================================================
%  Copy-paste LaTeX snippets to align each subsection with the
%  Newtonian (6.1.1) six-beat template.
%  Replace the \ref{...} labels with your actual labels.
%  Notation follows the manuscript: \eta_{\rm s}, \eta_{\rm p}, \lambda,
%  \sigma_{\rm exp}, \Delta x_{\max}.
% =====================================================================


% ---------------------------------------------------------------------
%  REUSABLE TEMPLATE  (paste at the top of any new subsection and fill)
% ---------------------------------------------------------------------
% (1) Setup + purpose.
%     We consider <geometry> with radius $a=1$ driven by <forcing> in a
%     <domain>. The fluid is <model> with true parameters <...>. This case
%     tests <purpose>.
% (2) FOM data.
%     The trajectory is simulated with the FOM until $t=<t_end>$ and sampled
%     at $N=<N>$ equally spaced points.
% (3) Noise ladder.
%     To mimic experimental data we add uncorrelated zero-mean Gaussian noise
%     with standard deviation <p>\% of the maximum displacement $\Delta x_{\max}$.
% (4) SAM-vs-FOM discrepancy decision.
%     The SAM evaluated at the ground-truth parameters <matches / does not match>
%     the FOM (Fig.~\ref{...}), so we <drop / retain> the model-bias term and
%     infer <parameter list> together with $\sigma_{\rm exp}$.
% (5) Priors.  The priors are listed in Table~\ref{...}.
% (6) Posterior.  Table~\ref{...} reports the posterior; <one-line takeaway>.


% ---------------------------------------------------------------------
%  6.1.2  Isolated viscoelastic  --  FIX beat 4 (replace existing sentence)
% ---------------------------------------------------------------------
As Fig.~\ref{fig:ve_iso_noise} shows, and as expected for this case, the SAM
using the ground-truth rheological parameters (red curve) closely matches the
FOM trajectories (black dots), so there is no systematic model discrepancy to
account for. We therefore drop the model-bias term for this case and infer the
three material parameters $\eta_{\rm s}$, $\eta_{\rm p}$, and $\lambda$ together
with the experimental-noise scale $\sigma_{\rm exp}$.


% ---------------------------------------------------------------------
%  6.1.3  Nonlinear viscoelastic  --  ADD beats 2--3 (up front)
% ---------------------------------------------------------------------
The FOM data are generated from $t=0$ to $t=0.5$ with the force removed at
$t=0.2$, sampled at $N=51$ equally spaced points. To mimic experimental data we
add uncorrelated zero-mean Gaussian noise with standard deviation $2\%$ of the
maximum displacement $\Delta x_{\max}$.

%  6.1.3  --  ADD beat 4 (the flipped decision)
In contrast to the isolated Newtonian and linear viscoelastic cases, the SAM
evaluated at the ground-truth parameters no longer matches the FOM trajectory in
this high-Weissenberg regime (Fig.~\ref{fig:nl_fom_sam}). A systematic model
discrepancy is therefore present, and we retain the model-bias term rather than
dropping it.


% ---------------------------------------------------------------------
%  6.2.1  Wall / Newtonian  --  ADD beats 3--4 (before the identifiability study)
% ---------------------------------------------------------------------
Before assessing identifiability, we confirm that the wall-corrected SAM at the
ground-truth parameters reproduces the FOM trajectory (Fig.~\ref{fig:wall_fom_sam}),
so no model-bias term is required for the bounded-domain forward model. As in the
isolated cases, we sample $N=30$ equally spaced points and add zero-mean Gaussian
noise with standard deviation $2\%$ of $\Delta x_{\max}$.

%  6.2.1  --  optional framing line for the add-on studies
The remaining analyses in this subsection go beyond the baseline recipe: they
quantify \emph{when} the two parameters become jointly identifiable
(curvature-to-noise threshold) and \emph{how} the inference degrades when the
wall correction is omitted from the forward model.


% ---------------------------------------------------------------------
%  6.2.2  Wall / viscoelastic  --  ADD beats 2--4
% ---------------------------------------------------------------------
Each of the three creep-recovery trajectories is sampled at $N=51$ equally spaced
points between $t=0$ and $t=0.5$, and we add zero-mean Gaussian noise with standard
deviation $2\%$ of $\Delta x_{\max}$. The SAM evaluated at the ground-truth
parameters matches the FOM response (Fig.~\ref{fig:wall_ve_fom_sam}), so no
model-bias term is required and we infer $\eta_{\rm s}$, $\eta_{\rm p}$, $\lambda$,
and the wall separation $\delta_0$ together with $\sigma_{\rm exp}$.


% ---------------------------------------------------------------------
%  6.3.1  Homogeneous material  --  ADD beat 4 framing
% ---------------------------------------------------------------------
The model-adequacy check in this setting is the periodic-box correction:
Table~\ref{tab:homog} contrasts the uncorrected inference with the
Hasimoto-corrected estimate, and the residual box-size dependence plays the role
of the SAM-vs-FOM discrepancy decision made in the earlier cases.


% ---------------------------------------------------------------------
%  6.3.2  Heterogeneous material  --  ADD priors/noise callout
% ---------------------------------------------------------------------
As in the homogeneous case, each trajectory is sampled at $N=51$ points between
$t=0$ and $t=0.5$ with $2\%$ Gaussian noise, using the priors in
Table~\ref{tab:pp_priors}. Here the model-adequacy question is different: rather
than checking a single trajectory, we compare the single-probe local estimate
against the bulk modulus (Fig.~\ref{fig:hetero_hist}).
